\chapter{合卷二：白银、军费与地方社会}
\label{ch:ming-synthesis-silver}

十六、十七世纪的白银常被写成一股从海外涌入、然后使明朝兴衰的单一潮水。这样的说法过于整齐。银既是田赋折纳、盐课、军饷、借贷和市场结算的媒介，也是必须被开采、运输、验色、兑换、借出和催收的实物。它连接中国内地、东南港口、日本、美洲与东南亚的贸易，却不会以同一速度、同一种成色或同一种风险抵达每个县份。

\section{一条鞭不是一条直线}

一条鞭将若干税役折并，扩大以田亩计征、以银两结算的范围。它减轻了某些旧有差役的碎片化，也可能使缺银家庭更依赖市场、地主或经手人。实物、劳役、临时加派和地方变通没有因此消失。研究改革时，法令、赋役全书、清丈册、契约和地方志必须互相校验：前者说明制度意图，后者才可能让我们看见一处地方怎样付款、欠缴、逃避或重新分摊。

\section{军费如何穿过社会}

宁夏、朝鲜、东南海防和辽东不是互不相干的花销。每一处都要通过仓储、漕运、商人垫解、盐课、州县催征和军户家庭获得资源。北京可以在账面上把它们列为不同项目；对沿途的役夫、船户、商人和纳税人而言，它们会在同一季节争夺劳力、粮食和信用。故而“财政枯竭”应当是要被解释的结果，不是可以跳过征解过程的原因。

\section{市场扩大不等于人人富裕}

江南市镇、矿业、手工业、海贸和长途贩运让更多人进入货币交换，也使价格、借贷和区域供应更能影响家庭。宗族、会馆、寺观、书院、牙行和地方公议并非经济之外的装饰；它们提供信用、救济、名望和信息，同时也可能排除没有关系或没有资产的人。材料对富户、士人和官员的保存往往较好，对佃户、妇女、雇工、船户和逃难者的声音较少。社会史必须承认这种沉默，而不能将保存下来的精英文书当作全体经验。

\section{危机不是一个总数}

晚明的旱灾、疫病、河患与战争，在不同区域先后穿过这张银、粮和人力的网络。一份灾报、一道蠲免令或一项饷额，只能证明它进入行政程序；不能自动证明救济到达、负担平均或反抗必然发生。末年的崩解因此应被看作多条反馈回路：战事提高筹饷压力，征解改变地方生计，灾害损伤道路和收成，信息迟滞又使调度更难。白银是其中重要的媒介，却不是替代战争、环境、政治选择和地方社会的万能答案。
